\begin{proposition}[Constant twists on an elliptic ruled surface]\hypertarget{tag-01AR}{}\label{cdp:st:A.4.2}
\StatementRecord{cdp:st:A.4.2}{proved}{01AR}{the definitions and constructions stated in the proof}

For an elliptic curve $E$, the minimal ruled surface $S_0=E\times\PP^1$ and the constant family $M_t=\OO_{S_0}$ satisfy $H^0(S_0,\Omega^1_{S_0}\otimes M_t)\ne0$ for every parameter. The ruling has genus one. These data give a counterexample to the genus-at-least-two inference in \cite[Proposition~7.3 and Step~1, p.~693]{CDP20}, where the line-bundle family is permitted to be constant.
\end{proposition}


\begin{proof}
Take $S_0=E\times\PP^1$ for an elliptic curve $E$ and take $M_t=\OO_{S_0}$ for every parameter. Every rational curve on $S_0$ maps constantly to $E$ by Riemann--Hurwitz, so it is a ruling fibre and has self-intersection zero. Thus $S_0$ has no exceptional curve of self-intersection $-1$ and is minimal. The product cotangent decomposition gives
\[
 h^q(S_0,\Omega^1_{S_0})=(1,2,1),\qquad q=0,1,2.
\]
Indeed the $E$ cotangent summand contributes $(1,1,0)$, while the $\PP^1$ cotangent summand contributes $(0,1,1)$ by the K\"unneth formula and $h^1(\PP^1,\OO(-2))=1$. A nonzero invariant form on $E$ provides the asserted sections for every parameter, whereas the ruling has genus one.

\end{proof}